% !TeX program = pdflatex
%
% can_micro.tex — Leitura de CAN/J1939 pelo Microprocessador Fractal
%
% Ponte entre o microprocessador fractal e um barramento CAN de veículo pesado.
%
% Autor: Aarão Melo Lopes
%

\documentclass[11pt,a4paper]{article}

\usepackage[utf8]{inputenc}
\usepackage[T1]{fontenc}
\usepackage[brazil]{babel}
\usepackage{amsmath,amssymb,amsthm}
\usepackage[margin=2.4cm]{geometry}
\usepackage{booktabs}
\usepackage{array}
\usepackage{tabularx}
\usepackage{listings}
\usepackage[hidelinks]{hyperref}

\lstset{
  basicstyle=\ttfamily\small,
  breaklines=true,
  columns=fullflexible,
  keepspaces=true,
  frame=single,
  framerule=0.4pt,
  xleftmargin=1em,
  xrightmargin=1em,
  aboveskip=0.8em,
  belowskip=0.8em
}

\newtheorem{proposicao}{Proposição}
\newtheorem{observacao}{Observação}

\title{\textbf{Microprocessador Fractal --- leitura de CAN/J1939}\\[0.4em]
\large do barramento físico ao estado do veículo}

\author{Aarão Melo Lopes}
\date{}

\begin{document}

\maketitle

\begin{abstract}

Este documento estabelece uma ponte entre o Microprocessador Fractal e a leitura
de um barramento CAN utilizado em veículos pesados. O objetivo não é afirmar que
CAN seja literalmente a operação $\oplus$ do corpo, mas mostrar como o mesmo
microprocessador pode ser empregado para receber, armazenar, decodificar e
transformar os frames CAN em grandezas físicas.

A cadeia é:

\[
\boxed{
\text{CAN físico}
\rightarrow
\text{frame}
\rightarrow
\text{identificador}
\rightarrow
\text{PGN}
\rightarrow
\text{SPN}
\rightarrow
\text{bytes}
\rightarrow
\text{escala}
\rightarrow
\text{estado físico}
}
\]

No vocabulário do Microprocessador Fractal, a entrada é uma sequência discreta
que chega pelo I/O, o identificador funciona como assinatura, o controle seleciona
o interpretador, a memória conserva o fluxo e a tríade $\oplus,\otimes,\int$
realiza as transformações necessárias.

Assim, ler CAN não exige introduzir um novo princípio computacional: exige apenas
uma nova assinatura sobre o mesmo corpo de cálculo.

\end{abstract}


% ─────────────────────────────────────────────────────────────────────────────
\section{O problema}

Um caminhão moderno possui diversos controladores eletrônicos:

\begin{center}
\begin{tabular}{c}
ECU do motor
\\
$\downarrow$
\\
transmissão --- ABS/EBS --- painel --- retarder --- telemetria
\\
$\uparrow$
\\
CAN
\end{tabular}
\end{center}

Esses controladores compartilham um meio físico. Cada participante pode produzir
mensagens e os demais podem observar o tráfego.

O problema do leitor não é inicialmente ``entender o caminhão''. O problema é
receber uma sequência de bits e descobrir a transformação que leva essa sequência
a uma grandeza com significado físico.

O Microprocessador Fractal já possui exatamente os ofícios necessários:

\[
\text{I/O}
\rightarrow
\text{controle}
\rightarrow
\text{memória}
\rightarrow
\text{núcleo}.
\]

A leitura CAN é, portanto, um caso de entrada do mundo físico para o corpo.


% ─────────────────────────────────────────────────────────────────────────────
\section{CAN como entrada física}

O CAN utiliza dois condutores diferenciais:

\[
\mathrm{CANH},\qquad \mathrm{CANL}.
\]

O transceptor converte a diferença elétrica entre esses condutores em bits digitais.
O microprocessador não precisa operar diretamente sobre a tensão do barramento:

\[
\boxed{
\text{CANH/CANL}
\rightarrow
\text{transceptor}
\rightarrow
\text{bits}
}
\]

A partir desse ponto, o barramento deixa de ser uma grandeza analógica para o
processador e passa a ser uma sequência discreta.

Um frame pode ser representado abstratamente por

\[
F=(I,D,L),
\]

onde

\begin{itemize}
\item $I$ é o identificador;
\item $D$ é a sequência de dados;
\item $L$ é o comprimento dos dados.
\end{itemize}

Para o microprocessador:

\[
\boxed{
F=(I,d_0,d_1,\ldots,d_{L-1})
}
\]

é simplesmente uma palavra de entrada.


% ─────────────────────────────────────────────────────────────────────────────
\section{A primeira assinatura --- o identificador}

O identificador CAN não é ainda a grandeza física. Ele funciona como uma assinatura
da mensagem.

O leitor recebe:

\[
I \quad+\quad D
\]

e utiliza $I$ para selecionar a interpretação apropriada.

No caso de protocolos de veículos pesados, como J1939, identificadores estendidos
contêm campos que permitem determinar a classe da mensagem e, consequentemente,
o seu PGN.

A arquitetura pode ser representada como:

\[
\boxed{
I
\rightarrow
\mathrm{decode}
\rightarrow
\mathrm{PGN}
\rightarrow
\mathrm{decoder}
}
\]

Isso corresponde diretamente ao controle do Microprocessador Fractal:

\[
\text{PC}
\rightarrow
\text{opcode}
\rightarrow
\text{operação}.
\]

A diferença é que, no CAN, a ``operação'' não é ADD ou MUL: é o interpretador
associado àquela assinatura de mensagem.


% ─────────────────────────────────────────────────────────────────────────────
\section{PGN e SPN como assinaturas}

No J1939, o PGN identifica um grupo de parâmetros transmitidos juntos.

Dentro desse grupo aparecem os parâmetros individuais, os SPNs.

A leitura pode ser vista como uma árvore:

\[
\boxed{
\text{frame}
\rightarrow
\text{PGN}
\rightarrow
\text{SPN}
\rightarrow
\text{campo de bits}
}
\]

O mesmo fluxo de bytes pode, portanto, produzir interpretações diferentes conforme
a assinatura que o acompanha.

Isso é precisamente o princípio usado no Microprocessador Fractal:

\[
\text{mesmo corpo}
+
\text{assinatura}
=
\text{ofício}.
\]

O CAN fornece o fluxo; o identificador seleciona a régua; o decoder realiza a leitura.


% ─────────────────────────────────────────────────────────────────────────────
\section{O dado ainda não é uma grandeza física}

Considere uma mensagem abstrata:

\begin{lstlisting}
ID   = 18F00401
DATA = 12 34 56 78 9A BC DE F0
\end{lstlisting}

Os bytes

\[
12,\;34,\;56,\;78,\;9A,\;BC,\;DE,\;F0
\]

não são automaticamente ``RPM'', ``temperatura'' ou ``torque''.

Primeiro é necessário selecionar um campo de bits.

Se um parâmetro ocupa os bits $r,\ldots,s$, define-se:

\[
R =
\sum_{k=r}^{s}
b_k2^{k-r}.
\]

Em uma implementação binária convencional, essa extração pode ser realizada por
máscara e deslocamento:

\[
R=(D\mathbin{\&}M)\gg r.
\]

O Microprocessador Fractal interpreta isso como uma operação de I/O: o fluxo bruto
é contraído para o estado relevante.


% ─────────────────────────────────────────────────────────────────────────────
\section{A transformação para unidade física}

Depois da extração, o valor bruto $R$ ainda pode não ser a grandeza física.

A conversão típica possui a forma

\[
x=S R+O,
\]

onde

\begin{itemize}
\item $R$ é o valor bruto;
\item $S$ é o fator de escala;
\item $O$ é o offset;
\item $x$ é a grandeza física.
\end{itemize}

Portanto:

\[
\boxed{
\text{bits}
\rightarrow
R
\rightarrow
SR+O
\rightarrow
x
}
\]

A multiplicação por $S$ e a soma de $O$ pertencem diretamente ao núcleo:

\[
R
\xrightarrow{\otimes S}
SR
\xrightarrow{\oplus O}
SR+O.
\]

Aqui a tríade aparece de maneira operacional:

\[
\boxed{
\oplus \quad \otimes
}
\]

são suficientes para realizar a transformação linear.


% ─────────────────────────────────────────────────────────────────────────────
\section{A leitura completa}

A leitura de um parâmetro CAN pode então ser escrita como a composição

\[
\mathcal{R}
=
\mathcal{S}
\circ
\mathcal{E}
\circ
\mathcal{D}
\circ
\mathcal{I},
\]

onde:

\begin{center}
\begin{tabular}{ll}
$\mathcal{I}$ & recebe o frame;\\
$\mathcal{D}$ & decodifica a assinatura;\\
$\mathcal{E}$ & extrai os bits;\\
$\mathcal{S}$ & aplica escala e offset.
\end{tabular}
\end{center}

Explicitamente:

\[
F
\longmapsto
(I,D)
\longmapsto
(PGN,SPN)
\longmapsto
R
\longmapsto
SR+O.
\]

O resultado final é um estado físico:

\[
\boxed{x=SR+O}.
\]


% ─────────────────────────────────────────────────────────────────────────────
\section{O Microprocessador Fractal como leitor CAN}

A arquitetura completa fica:

\begin{center}
\begin{tabular}{c}
CANH/CANL
\\
$\downarrow$
\\
transceptor
\\
$\downarrow$
\\
frame $(I,D)$
\\
$\downarrow$
\\
controle / decode
\\
$\downarrow$
\\
PGN / SPN
\\
$\downarrow$
\\
I/O --- extração de bits
\\
$\downarrow$
\\
$\oplus,\otimes$
\\
$\downarrow$
\\
grandeza física
\\
$\downarrow$
\\
disco / estado
\end{tabular}
\end{center}

Cada bloco já possui um correspondente na arquitetura:

\begin{center}
\begin{tabular}{lll}
\toprule
CAN & Microprocessador & função\\
\midrule
frame & entrada & sequência discreta\\
identifier & assinatura & seleção do interpretador\\
PGN/SPN & controle & classificação\\
bit field & I/O & extração\\
scale/offset & núcleo & $\otimes,\oplus$\\
log/exp, se necessário & núcleo analógico & $\int$\\
registro do frame & disco & estado\\
\bottomrule
\end{tabular}
\end{center}


% ─────────────────────────────────────────────────────────────────────────────
\section{O barramento não é literalmente o $\oplus$}

Há uma distinção importante.

Não se afirma:

\[
\mathrm{CAN}=\oplus.
\]

O CAN possui um mecanismo específico de arbitragem e transmissão diferencial.
O $\oplus$ do Microprocessador Fractal é a operação algébrica de soma.

A correspondência correta é:

\[
\boxed{
\text{CAN}
\sim
\text{meio compartilhado distribuído}
}
\]

enquanto

\[
\boxed{
\oplus
=
\text{operação de soma do núcleo}.
}
\]

A relação entre ambos está no fato de que múltiplos participantes utilizam um
mesmo meio, mas os mecanismos físicos e algébricos não devem ser identificados.

Essa separação protege a arquitetura contra uma equivalência indevida.


% ─────────────────────────────────────────────────────────────────────────────
\section{A memória --- o fluxo CAN como disco}

O barramento produz uma sequência temporal:

\[
F_0,F_1,F_2,\ldots,F_t.
\]

O processador pode tratá-la como um arquivo:

\[
\mathcal{D}
=
(F_0,F_1,\ldots,F_t).
\]

O disco, no vocabulário do corpo, torna-se o histórico do barramento.

Uma janela temporal pode então ser contraída:

\[
\mathcal{D}
\rightarrow
\mathcal{D}_{\mathrm{PGN}}
\rightarrow
\mathcal{D}_{\mathrm{SPN}}
\rightarrow
\mathcal{D}_{x}.
\]

O mesmo princípio de contração utilizado pela memória do Microprocessador Fractal
aparece agora como filtragem e interpretação do fluxo.

O arquivo cresce com o tráfego; a palavra de trabalho necessária para interpretar
cada frame não precisa crescer proporcionalmente ao tamanho do histórico.


% ─────────────────────────────────────────────────────────────────────────────
\section{O tempo --- CAN como sequência de estados}

O barramento não fornece apenas valores; fornece valores no tempo.

Assim:

\[
x_0,x_1,x_2,\ldots
\]

constitui uma trajetória do veículo.

O leitor pode construir:

\[
\Delta x_t=x_t-x_{t-1},
\]

ou uma taxa aproximada:

\[
\dot{x}_t
\approx
\frac{x_t-x_{t-1}}{t_t-t_{t-1}}.
\]

Novamente, apenas as operações básicas aparecem:

\[
\ominus,\quad \otimes,\quad \otimes^{-1}.
\]

A sequência CAN transforma-se, portanto, em uma dinâmica discreta:

\[
x_{t+1}=F(x_t,u_t).
\]

O veículo passa a ser observado como uma sequência de estados.


% ─────────────────────────────────────────────────────────────────────────────
\section{Um leitor mínimo}

Um leitor conceitual pode ser escrito como:

\begin{lstlisting}
while (can_read(&frame)) {

    id   = frame.id;
    data = frame.data;

    pgn = decode_pgn(id);

    if (pgn == PGN_ALVO) {

        raw = extract_bits(data);

        value = raw * SCALE + OFFSET;

        store(value);
    }
}
\end{lstlisting}

A estrutura corresponde diretamente ao corpo:

\begin{lstlisting}
READ
  |
DECODE
  |
EXTRACT
  |
MUL
  |
ADD
  |
STORE
\end{lstlisting}

Não aparece um novo computador. Aparece uma nova aplicação do mesmo computador.


% ─────────────────────────────────────────────────────────────────────────────
\section{A assinatura como régua}

O ponto central desta integração pode ser escrito:

\[
\boxed{
\text{frame bruto}
+
\text{assinatura}
\longrightarrow
\text{grandeza}.
}
\]

A assinatura determina:

\[
\Sigma_{\mathrm{CAN}}
=
(\mathrm{ID},\mathrm{PGN},\mathrm{SPN},
\mathrm{posição},\mathrm{comprimento},
\mathrm{escala},\mathrm{offset}).
\]

Ela funciona como uma régua.

A mesma sequência de bits pode ser lida como diferentes objetos dependendo da
régua aplicada. O corpo computacional, entretanto, permanece:

\[
\oplus,\qquad\otimes,\qquad\text{controle},\qquad\text{estado}.
\]

Muda a assinatura; não muda o núcleo.


% ─────────────────────────────────────────────────────────────────────────────
\section{A ponte completa}

A integração pode finalmente ser resumida por:

\[
\boxed{
\begin{array}{ccccccccc}
\text{veículo}
&\rightarrow&
\text{CAN}
&\rightarrow&
\text{frame}
&\rightarrow&
\text{PGN/SPN}
&\rightarrow&
\text{estado físico}
\\
&&&&
\downarrow
&&
\downarrow
&&
\downarrow
\\
&&&
\text{I/O}
&
\rightarrow
&
\text{controle}
&
\rightarrow
&
\oplus,\otimes
\end{array}
}
\]

O barramento fornece a sequência.

A assinatura determina como lê-la.

O controle seleciona o caminho.

O I/O extrai o campo.

O núcleo transforma o valor.

O disco conserva o estado.

Assim, o Microprocessador Fractal pode ser utilizado como uma máquina de leitura de
um barramento CAN sem introduzir uma arquitetura externa ao corpo:

\[
\boxed{
\text{CAN}
\rightarrow
\text{I/O}
\rightarrow
\text{controle}
\rightarrow
\text{tríade}
\rightarrow
\text{estado}.
}
\]


% ─────────────────────────────────────────────────────────────────────────────
\section{Conclusão}

O CAN fornece um caso físico particularmente claro para a arquitetura do
Microprocessador Fractal.

Não porque CAN seja ``a mesma coisa'' que o corpo, mas porque o caminho de uma
grandeza física até uma decisão computacional apresenta exatamente a decomposição
que o corpo procura:

\[
\text{meio}
\rightarrow
\text{assinatura}
\rightarrow
\text{decodificação}
\rightarrow
\text{transformação}
\rightarrow
\text{estado}.
\]

Em um caminhão, uma rotação do motor nasce como fenômeno físico, é medida por um
sensor, convertida por uma ECU em dados digitais, transportada pelo CAN, identificada
por sua assinatura, extraída de um campo de bits e finalmente transformada em um
número físico.

No Microprocessador Fractal, esse caminho pode ser lido como:

\[
\boxed{
\text{I/O}
\rightarrow
\text{controle}
\rightarrow
\oplus,\otimes
\rightarrow
\text{memória}.
}
\]

O barramento externo torna-se, assim, uma nova fronteira de realização do mesmo
princípio: não uma nova máquina, mas o mesmo corpo recebendo uma nova régua.

\end{document}
